% ---------- Color Commands ----------
\DeclareRobustCommand{\yellow}[1]{\begingroup\color{Yellow}\ignorespaces#1\endgroup\ignorespacesafterend}
\DeclareRobustCommand{\red}[1]{\begingroup\color{Red}\ignorespaces#1\endgroup\ignorespacesafterend}
\DeclareRobustCommand{\blue}[1]{\begingroup\color{Blue}\ignorespaces#1\endgroup\ignorespacesafterend}
\DeclareRobustCommand{\green}[1]{\begingroup\color{Green}\ignorespaces#1\endgroup\ignorespacesafterend}
\DeclareRobustCommand{\gray}[1]{\begingroup\color{Gray}\ignorespaces#1\endgroup\ignorespacesafterend}
\newenvironment{redblock}{\begingroup\color{Red}\ignorespaces}{\endgroup\ignorespacesafterend}

\newenvironment{stenosisillustration}{%
    \begin{center}
    \begin{minipage}{0.92\linewidth}
    \hrule
    \smallskip
    \small
    \textbf{Illustration from the idealized stenosis study.}\par\smallskip
}{%
    \smallskip
    \hrule
    \end{minipage}
    \end{center}
}


% ---------- Define custom colors ----------
\definecolor{bg}{rgb}{0.95,0.95,0.92}
\definecolor{codeblue}{rgb}{0.25,0.5,0.75}
\definecolor{keyword}{rgb}{0.5,0,0.5}
\definecolor{string}{rgb}{0,0.5,0}

% ---------- Define the style for IPython ----------
\lstdefinelanguage{ipython}{
  basicstyle=\ttfamily\footnotesize\color{black},
  keywordstyle=\color{keyword}\bfseries,
  stringstyle=\color{string},
  backgroundcolor=\color{bg},
  morekeywords={In,Out},
  frame=lines,
  numbers=left,
  numbersep=5pt
}

% ---------- Define the style for Julia ----------
\lstdefinelanguage{julia}{
  keywords={struct, abstract, primitive, typealias, bitstype, mutable, module, baremodule, import, importall, using, export, if, else, elseif, for, while, break, continue, function, return, type, global, local, const, let, do, try, catch, finally, end, quote, macro, where},
  basicstyle=\ttfamily\footnotesize\color{black},
  keywordstyle=\color{keyword}\bfseries,
  stringstyle=\color{string},
  backgroundcolor=\color{bg},
  frame=lines,
  numbers=left,
  numbersep=8pt,
  inputencoding=utf8,
  extendedchars=true,
  escapeinside={(*}{*)},
  literate={~}{{$\sim$}}1
           {₁}{{\textsubscript{1}}}1
           {₂}{{\textsubscript{2}}}1
           {₃}{{\textsubscript{3}}}1
           {ᵀ}{{$^{\mathrm{T}}$}}1
           {⋅}{{$\cdot$}}1
           {∈}{{$\in$}}1
           {ℝ}{{$\mathbb{R}$}}1
           {→}{{$\to$}}1
           {π}{{$\pi$}}1
           {≤}{{$\leq$}}1
           {≥}{{$\geq$}}1
           {≈}{{$\approx$}}1
           {ϵ}{{$\epsilon$}}1
           {κ}{{$\kappa$}}1
           {Δ}{{$\Delta$}}1
           {λ}{{$\lambda$}}1
}

% ---------- Define the style for Mathematica ----------
\lstdefinelanguage{Mathematica}{
  keywords={Module,With,Block,If,Then,Else,Which,Switch,For,While,Return,Do,Table,Plot,Map,Apply,Function,Set,SetDelayed,Clear,Quit,NonlinearModelFit,FindMinimum,FindMaximum,FindRoot,NIntegrate,Integrate,Simplify,FullSimplify,DSolve,RSolve,NSolve,NDSolve,Limit,Series,Assuming,Expand,Factor,TableForm,MatrixForm,Part,Length,Dimensions,Transpose,MapThread,MapAt,Flatten,Thread,Join,Outer,ConstantArray,Riffle,ArrayPlot,Plot3D,ContourPlot,ParametricPlot,MatrixPlot,Graphics,Graphics3D,Manipulate,Evaluate,Sin,Cos,Exp,Log,Sqrt},
  keywordstyle=\color{keyword}\bfseries,
  sensitive=true,
  morecomment=[l]{(*}, morecomment=[r]{*)},
  commentstyle=\itshape\color{gray},
  morestring=[b]",
  stringstyle=\color{string},
  basicstyle=\ttfamily\footnotesize\color{black},
  backgroundcolor=\color{bg},
  frame=lines,
  numbers=left,
  numbersep=5pt,
  showstringspaces=false,
  escapeinside={(*}{*)}
}

% ---------- Optionally define a style for Mathematica code cells ----------
\lstdefinestyle{mathematicaStyle}{
  language=Mathematica,
  frame=lines,
  backgroundcolor=\color{bg},
  basicstyle=\ttfamily\footnotesize,
  keywordstyle=\color{keyword}\bfseries,
  commentstyle=\itshape\color{gray},
  stringstyle=\color{string},
  numbers=left,
  numbersep=5pt,
  showstringspaces=false
}

% ============================================================
% =================  FIGURE STYLES (TikZ)  ===================
% ============================================================
% Shared palette and styles used by figures/static/static/tikz/*.tex.
% Keep these consistent across this report for a unified visual
% language: blood = warm red; walls = neutral gray; mathematical
% annotations = dark blue; map/context artifacts = teal.
% Effective 6-color palette: mathblue (primary annotation), accentTeal
% (map / context secondary), bloodred (blood / flow / severity),
% wallgray (structure), plus light fills for the three primary hues.
% stenosisOrange and accentGold are retained as named aliases for
% backward compatibility with figure source files but are now muted
% members of the bloodred and accentTeal families (no new hues).
\definecolor{bloodred}{HTML}{B23A3A}
\definecolor{bloodredlight}{HTML}{F2C7C2}
\definecolor{wallgray}{HTML}{6E6E6E}
\definecolor{wallgraylight}{HTML}{D8D5D2}
\definecolor{mathblue}{HTML}{1F3A5F}
\definecolor{mathblueLite}{HTML}{C9D6E5}
\definecolor{accentTeal}{HTML}{217A7A}
\definecolor{accentTealLite}{HTML}{C7E4E1}
% Retired hues, remapped to within-family muted variants so all
% legacy `stenosisOrange` / `accentGold` call sites still resolve
% without introducing a third or fourth chromatic family.
\colorlet{stenosisOrange}{bloodred!72!black}
\colorlet{accentGold}{accentTeal!72!black}

\tikzset{
  % container styles
  fig/.style={
    every node/.append style={font=\footnotesize},
    >=Latex, line cap=round, line join=round, thick
  },
  vesselWall/.style={draw=wallgray, line width=0.9pt, fill=wallgraylight},
  lumen/.style    ={draw=bloodred!70!black, line width=0.6pt, fill=bloodredlight},
  centerline/.style={dashed, gray!70, line width=0.5pt},
  axisArrow/.style ={->, thick, mathblue, line width=0.7pt},
  flowArrow/.style ={->, thick, bloodred!80!black, line width=0.9pt},
  labelM/.style    ={font=\footnotesize, mathblue},
  labelB/.style    ={font=\footnotesize, bloodred!70!black},
  labelW/.style    ={font=\footnotesize, wallgray!50!black},
  % --- Shared figure aliases -------------------------------------
  % Flow / primary action arrows (vessel inflow, dominant edges).
  flowArrowMain/.style ={->, line cap=round, line join=round,
                         bloodred!80!black, line width=1.0pt},
  % Schematic flowchart edge (block-to-block).
  schemEdge/.style ={->, line cap=round, line join=round,
                     mathblue!75!black, line width=0.9pt},
  % Measurement / axis / dimension arrows.
  measureArrow/.style ={->, mathblue, line width=0.7pt},
  measureBiArrow/.style ={<->, mathblue, line width=0.7pt},
  % Light helper / reference / dashed guide.
  helperLine/.style ={dashed, gray!55, line width=0.5pt},
  % Unified label background variants (white halo, opacity 0.92).
  tagLabel/.style ={font=\footnotesize, fill=white, fill opacity=0.92,
                    text opacity=1, inner sep=1.2pt},
  tagLabelM/.style ={font=\footnotesize, text=mathblue, fill=white,
                     fill opacity=0.92, text opacity=1, inner sep=1.2pt},
  tagLabelB/.style ={font=\footnotesize, text=bloodred!70!black, fill=white,
                     fill opacity=0.92, text opacity=1, inner sep=1.2pt},
  tagLabelT/.style ={font=\footnotesize, text=accentTeal!70!black, fill=white,
                     fill opacity=0.92, text opacity=1, inner sep=1.2pt},
  panelTitle/.style ={font=\small\bfseries, text=mathblue,
                      inner sep=2pt},
  % boxed schematic block (for flowcharts / hierarchy)
  schemBlock/.style={
    draw=mathblue, line width=0.8pt, rounded corners=2pt,
    fill=mathblueLite!45, inner sep=4pt, align=center, font=\footnotesize,
    minimum height=9mm, minimum width=26mm
  },
  schemOperator/.style={
    draw=accentTeal, line width=1pt, rounded corners=3pt,
    fill=accentTealLite, inner sep=5pt, align=center, font=\footnotesize,
    minimum height=10mm, minimum width=26mm
  },
  schemData/.style={
    draw=accentGold, line width=0.8pt,
    fill=accentTealLite!55, inner sep=4pt, align=center, font=\footnotesize,
    minimum height=9mm, minimum width=26mm,
    chamfered rectangle, chamfered rectangle xsep=4pt
  },
  schemArrow/.style={->, line width=0.9pt, mathblue!80!black},
  % small markers
  pulseDot/.style={circle, fill=bloodred, inner sep=1.2pt},
  ptDot/.style={circle, fill=mathblue, inner sep=1pt}
}

% --- Shared pgfplots defaults for figure panels -------------------
% Applied automatically to every axis environment in figure files.
% Light gray grid, footnotesize axis labels, small bold titles,
% mathblue axis-line color, ticks outside.
\pgfplotsset{
  every axis/.append style={
    axis line style={mathblue!55!black, line width=0.45pt},
    tick style={mathblue!55!black, line width=0.45pt},
    grid style={gray!22, line width=0.3pt},
    tick align=outside,
    tick label style={font=\scriptsize, mathblue!85!black},
    label style={font=\footnotesize, mathblue},
    title style={font=\small\bfseries, mathblue, yshift=-1pt},
    legend style={font=\scriptsize, fill=white, fill opacity=0.92,
                  text opacity=1, draw=mathblue!30, line width=0.4pt,
                  inner sep=2pt, rounded corners=1pt},
  },
  sevpanel/.style={width=5.75cm, height=3.35cm,
    tick label style={font=\scriptsize},
    label style={font=\scriptsize},
    xlabel style={font=\scriptsize},
    ylabel style={font=\scriptsize},
    title style={font=\footnotesize\bfseries, yshift=-1pt},
    xlabel={$t/T_h$}, xmin=0, xmax=1,
    xtick={0,0.5,1},
    major tick length=2pt,
    grid=both, grid style={dashed, gray!25},
    every axis plot/.append style={line width=0.9pt}},
  signalpanel/.style={width=7.1cm, height=4.55cm,
    tick label style={font=\scriptsize},
    label style={font=\scriptsize},
    xlabel style={font=\scriptsize},
    ylabel style={font=\scriptsize},
    title style={font=\footnotesize\bfseries, yshift=-1pt},
    xlabel={$t/T_h$}, xmin=0, xmax=1,
    xtick={0,0.5,1},
    major tick length=2pt,
    grid=both, grid style={dashed, gray!25},
    every axis plot/.append style={line width=0.9pt}},
}

% Compact appendix tables should use this wrapper instead of local
% \arraystretch overrides in body files.
\newenvironment{compactacronymtables}
  {\begingroup\footnotesize\renewcommand{\arraystretch}{1.08}}
  {\endgroup}

% ---------- Figure-input convenience macros ----------
% \figtikz{name} -> \input{figures/static/static/tikz/name.tex}
\newcommand{\figtikz}[1]{\input{figures/static/static/tikz/#1.tex}}
